\section{DevDay Surface}
\sectionkicker{PRODUCT SURFACE}
\begin{wideflow}
{\Large\bfseries DevDayの製品面――コンピューター操作API、常駐エージェントdots、発表の束ね}\par
\smallskip
{\color{SurveyMuted}変更履歴の文言の範囲で読む。何が出荷されたかの記録であり、信頼性の証拠ではない。}\par
\end{wideflow}

OpenAIのAgents APIに、ホスト型ブラウザーでのコンピューター操作機能が加わった\cite{agentsapi}。人が預かる計算機を、認可やサインインの始末を利用者側の手続きに任せて操るもので、APIとCodex、WorkのPro500やEnterpriseプランとの組み合わせで提供される。変更履歴の逐語が根拠であり、利用方法は読み取れるが、ここからブラウザーの安全性や信頼性を引き出してはならない。操作手順の解説文書は未読である。

dotsは、Astraの技術を土台とする常駐型エージェントという位置づけである\cite{dots}。クラウド上のコンピューターとブラウザー、4000超のプラグイン、読み取り専用の先回り動作、自動レビューが並ぶ。ProやBiz-PremiumとEnterprise・Eduのベータ提供の組み合わせも示される。製品説明の範囲として読み、測定された成果とは切り分ける。

DevDayの発表の束ねは、30番目の見出しにはしない\cite{devday}。Ultrafastプラン、Decisionsプレビュー、エージェントのコンピューター操作、Codexファミリー、プラグインやSpaceやPro500やマーケットプレイスが一束になった目録であり、支えとして小さく載せる。見出しに立てれば製品カタログの羅列になるだけである。

\begin{claimboundary}[資料と評価の境界]
変更履歴は、何が出荷されたかの記録であり、見守りや確からしさの証拠ではない。操作手順の解説文書は未読である。製品説明の文言を測定された成果と読み替えない。束ねは支えであり、見出しにはしない。
\end{claimboundary}
